\section{Limitations}
\label{sec:limits}

\begin{enumerate}\itemsep3pt

\item \textbf{The taste scale $\sigma_{\varepsilon}$ is calibrated, not
estimated.} $\psi_g$ and $\sigma_\varepsilon$ jointly hit $D^G_{ss}=0.05$, and one
level moment cannot separately identify a switching cost and a taste scale.
Section~\ref{sec:taste_id} characterises the locus: the steady state and the
policy ordering are invariant to $\sigma_\varepsilon$, but the adoption-channel
magnitudes are not, and pinning them needs an adoption-elasticity moment
(elasticity $1.3$ at our baseline). Estimating that elasticity directly, rather
than reading it off the micro literature, is the natural next step.

\item \textbf{The switching cost is booked as an import by accounting
necessity.} Under AMRS's $\mu_{ss}>1$ with capitalised profits, booking
$\psi_g D^{sw}$ as domestic demand creates a marginal dividend owned by nobody
and breaks asset-market clearing; we therefore book it externally
(Section~\ref{sec:bop}). This convention is what makes the adoption channel
\emph{contractionary} under the price shock, since the cost lands in the
periods household income is falling hardest (Section~\ref{sec:crisis}). The \emph{ordering}
of the policy results does not depend on it, it follows from the relative
price alone, but the sign of the adoption channel's output contribution under
a price shock does, and should be read as conditional on this booking.

\item \textbf{No modelled green technology.} The adopters' resource saving
rests on the domestic/imported interpretation of Section~\ref{sec:bop} rather
than on a green sector with its own capacity, cost curve and investment. We
have implemented the minimal version, a competitive, zero-profit installer
sector that books the switching cost as domestic value added rebated to
households, and it satisfies market clearing exactly and flips the sign of the
adoption channel (Section~\ref{sec:signmap}); what remains unmodelled is the
green energy \emph{supply} technology (capacity, cost curve, investment) and a
distribution of installer income more realistic than the lump-sum rebate.

\item \textbf{No within-country energy-share heterogeneity.}
\citet{Bayer2026}'s central distributional mechanism is that poor households
have higher energy shares; we have homothetic energy demand conditional on
the durable type, so we cannot speak to their distributional results.

\item \textbf{No investment margin.} Both \citet{Bayer2026} and
\citet{ARS2024} have richer supply sides. Its absence here likely makes the
transfer look more expansionary, and the supply shock more inflationary, than
it would be with capital to smooth the adjustment.

\end{enumerate}
